\documentclass[10pt, letterpaper]{article}

% Packages:
\usepackage[
    ignoreheadfoot, % set margins without considering header and footer
    top=2 cm, % seperation between body and page edge from the top
    bottom=2 cm, % seperation between body and page edge from the bottom
    left=2 cm, % seperation between body and page edge from the left
    right=2 cm, % seperation between body and page edge from the right
    footskip=1.0 cm, % seperation between body and footer
    % showframe % for debugging
]{geometry} % for adjusting page geometry
\usepackage{titlesec} % for customizing section titles
\usepackage{tabularx} % for making tables with fixed width columns
\usepackage{array} % tabularx requires this
\usepackage[dvipsnames]{xcolor} % for coloring text
\definecolor{primaryColor}{RGB}{0, 79, 144} % define primary color
\usepackage{enumitem} % for customizing lists
\usepackage{fontawesome5} % for using icons
\usepackage{amsmath} % for math
\usepackage[
    pdftitle={Pablo Morilla Cabello's CV},
    pdfauthor={Pablo Morilla Cabello},
    pdfcreator={LaTeX with Overleaf},
    colorlinks=true,
    urlcolor=primaryColor
]{hyperref} % for links, metadata and bookmarks
\usepackage[pscoord]{eso-pic} % for floating text on the page
\usepackage{calc} % for calculating lengths
\usepackage{bookmark} % for bookmarks
\usepackage{lastpage} % for getting the total number of pages
\usepackage{changepage} % for one column entries (adjustwidth environment)
\usepackage{paracol} % for two and three column entries
\usepackage{ifthen} % for conditional statements
\usepackage{needspace} % for avoiding page brake right after the section title
\usepackage{iftex} % check if engine is pdflatex, xetex or luatex

% Ensure that generate pdf is machine readable/ATS parsable:
\ifPDFTeX
    \input{glyphtounicode}
    \pdfgentounicode=1
    % \usepackage[T1]{fontenc} % this breaks sb2nov
    \usepackage[utf8]{inputenc}
    \usepackage{lmodern}
\fi



% Some settings:
\AtBeginEnvironment{adjustwidth}{\partopsep0pt} % remove space before adjustwidth environment
\pagestyle{empty} % no header or footer
\setcounter{secnumdepth}{0} % no section numbering
\setlength{\parindent}{0pt} % no indentation
\setlength{\topskip}{0pt} % no top skip
\setlength{\columnsep}{0cm} % set column seperation
\makeatletter
\let\ps@customFooterStyle\ps@plain % Copy the plain style to customFooterStyle
\patchcmd{\ps@customFooterStyle}{\thepage}{
    \color{gray}\textit{\small Pablo Morilla Cabello - Page \thepage{} of \pageref*{LastPage}}
}{}{} % replace number by desired string
\makeatother
\pagestyle{customFooterStyle}

\titleformat{\section}{\needspace{4\baselineskip}\bfseries\large}{}{0pt}{}[\vspace{1pt}\titlerule]

\titlespacing{\section}{
    % left space:
    -1pt
}{
    % top space:
    0.3 cm
}{
    % bottom space:
    0.2 cm
} % section title spacing

\renewcommand\labelitemi{$\circ$} % custom bullet points
\newenvironment{highlights}{
    \begin{itemize}[
        topsep=0.10 cm,
        parsep=0.10 cm,
        partopsep=0pt,
        itemsep=0pt,
        leftmargin=0.4 cm + 10pt
    ]
}{
    \end{itemize}
} % new environment for highlights

\newenvironment{highlightsforbulletentries}{
    \begin{itemize}[
        topsep=0.10 cm,
        parsep=0.10 cm,
        partopsep=0pt,
        itemsep=0pt,
        leftmargin=10pt
    ]
}{
    \end{itemize}
} % new environment for highlights for bullet entries


\newenvironment{onecolentry}{
    \begin{adjustwidth}{
        0.2 cm + 0.00001 cm
    }{
        0.2 cm + 0.00001 cm
    }
}{
    \end{adjustwidth}
} % new environment for one column entries

\newenvironment{twocolentry}[2][]{
    \onecolentry
    \def\secondColumn{#2}
    \setcolumnwidth{\fill, 5 cm}
    \begin{paracol}{2}
}{
    \switchcolumn \raggedleft \secondColumn
    \end{paracol}
    \endonecolentry
} % new environment for two column entries

\newenvironment{header}{
    \setlength{\topsep}{0pt}\par\kern\topsep\centering\linespread{1.5}
}{
    \par\kern\topsep
} % new environment for the header

\newcommand{\placelastupdatedtext}{% \placetextbox{<horizontal pos>}{<vertical pos>}{<stuff>}
  \AddToShipoutPictureFG*{% Add <stuff> to current page foreground
    \put(
        \LenToUnit{\paperwidth-2 cm-0.2 cm+0.05cm},
        \LenToUnit{\paperheight-1.0 cm}
    ){\vtop{{\null}\makebox[0pt][c]{
        \small\color{gray}\textit{Last updated in November 2026}\hspace{\widthof{Last updated in November 2026}}
    }}}%
  }%
}%

% save the original href command in a new command:
\let\hrefWithoutArrow\href

% new command for external links:
\renewcommand{\href}[2]{\hrefWithoutArrow{#1}{\ifthenelse{\equal{#2}{}}{ }{#2 }\raisebox{.15ex}{\footnotesize \faExternalLink*}}}


\begin{document}
    \newcommand{\AND}{\unskip
        \cleaders\copy\ANDbox\hskip\wd\ANDbox
        \ignorespaces
    }
    \newsavebox\ANDbox
    \sbox\ANDbox{}

    \placelastupdatedtext
    \begin{header}
        \textbf{\fontsize{24 pt}{24 pt}\selectfont Pablo Morilla Cabello}

        \vspace{0.3 cm}

        \normalsize
        \mbox{{\color{black}\footnotesize\faMapMarker*}\hspace*{0.13cm}Málaga, Spain}%
        \kern 0.25 cm%
        \AND%
        \kern 0.25 cm%
        \mbox{\hrefWithoutArrow{mailto:pablomorillacabello@gmail.com}{\color{black}{\footnotesize\faEnvelope[regular]}\hspace*{0.13cm}pablomorillacabello@gmail.com}}%
        \kern 0.25 cm%
        \AND%
        \kern 0.25 cm%
        \mbox{\hrefWithoutArrow{tel:+34 626 171 323}{\color{black}{\footnotesize\faPhone*}\hspace*{0.13cm}+34 626 171 323}}%
        \kern 0.25 cm%
        \AND%
        \kern 0.25 cm%
        \mbox{\hrefWithoutArrow{https://linkedin.com/in/PabloMorillaCabello}{\color{black}{\footnotesize\faLinkedinIn}\hspace*{0.13cm}PabloMorillaCabello}}%
        \kern 0.25 cm%
        \AND%
        \kern 0.25 cm%
        \mbox{\hrefWithoutArrow{https://github.com/PabloMorillaCabello}{\color{black}{\footnotesize\faGithub}\hspace*{0.13cm}PabloMorillaCabello}}%
        \AND%
        \kern 0.25 cm%
        \mbox{\hrefWithoutArrow{https://pablomorillacabello.github.io/}{\color{black}{\footnotesize\faGlobe[regular]}\hspace*{0.13cm}Portfolio}}%
        \kern 0.25 cm%
    \end{header}

    \vspace{0.3 cm - 0.3 cm}

    \section{Experience}

        \begin{twocolentry}{
        \textit{Málaga, Spain}

        \textit{August 2026 - Present}}
        \textbf{Development Engineer}
        \textit{\href{https://dhvtechnology.com/}{DHV Technology}}
        \end{twocolentry}

        \vspace{0.10 cm}
        \begin{onecolentry}
            \begin{highlights}
            \item Development engineering of space-grade solar power systems for small satellites and CubeSats.
            \item Applying industrial automation and robotics experience to the development and testing of space hardware.
            \end{highlights}
        \end{onecolentry}

        \vspace{0.2 cm}

        \begin{twocolentry}{
        \textit{Skövde, Sweden}

        \textit{April 2024 - June 2026}}
        \textbf{Automation Engineer}
        \textit{\href{https://elektroautomatik.se/}{Elektroautomatik}}
        \end{twocolentry}

        \vspace{0.10 cm}
        \begin{onecolentry}
            \begin{highlights}
            \item Programmed and commissioned industrial robot cells and PLC-based automation lines in multidisciplinary teams.
            \item Developed reusable programming standards to optimize workflows.
            \item Worked with RAPID, TIA Portal, TwinCAT and CODESYS, from offline programming to system integration.
            \item Delivered on-site commissioning for international companies such as Volvo Group, SCANIA, and Hitachi.
            \end{highlights}
        \end{onecolentry}

        \vspace{0.2 cm}

        \begin{twocolentry}{
        \textit{Skövde, Sweden}

        \textit{June 2023 - August 2023}}
        \textbf{Research Assistant}
        \textit{\href{https://assarinnovation.se/}{ASSAR Industrial Innovation Arena}}
      \end{twocolentry}

        \vspace{0.10 cm}
        \begin{onecolentry}
            \begin{highlights}
            \item Planned a summer project within a larger research group.
            \item Developed and integrated dedicated software into robotic and automation systems.
            \item Simulated a 3D industrial robotic system.
            \item Contributed to writing and presenting a published paper after the project.
            \end{highlights}
        \end{onecolentry}

    \section{Education}

        \begin{twocolentry}{


        \textit{August 2025 – June 2026}}
            \textbf{University West}

            \textit{MS in AI and Automation}
        \end{twocolentry}

        \vspace{0.10 cm}
        \begin{onecolentry}
            \begin{highlights}
            \item \href{https://www.hv.se/utbildning/program/master-i-ai-och-automation-heltid-campus-tamai/}{Master's WebPage}

                    \item \textbf{Coursework:} Machine Learning, DNN, CNN, RNN, Python, TensorFlow, pipeline design, OpenCV, computer vision, PLC programming, RAPID, ROS2, Docker.
            \end{highlights}
        \end{onecolentry}

        \vspace{0.20 cm}

        \begin{twocolentry}{


        \textit{August 2022 – June 2023}}
            \textbf{Högskolan i Skövde}

            \textit{BS in Industrial Engineering}
        \end{twocolentry}

        \vspace{0.10 cm}
        \begin{onecolentry}
            \begin{highlights}
            \item \href{https://www.his.se/en/education/engineering-and-technology/engineering-course-combinations/}{Bachelor's WebPage}

                    \item \textbf{Coursework:} Industrial Control Engineering, Industrial Robotics and Offline Programming, Logistics Simulation.
            \end{highlights}
        \end{onecolentry}

        \vspace{0.20 cm}

        \begin{twocolentry}{


        \textit{August 2019 – February 2024}}
            \textbf{University of Málaga}

            \textit{Degree in Robotics, Electronics, and Mechatronics}
        \end{twocolentry}

        \vspace{0.10 cm}
        \begin{onecolentry}
            \begin{highlights}
            \item \href{https://www.uma.es/grado-en-ingenieria-electronica-robotica-y-mecatronica}{Degree's WebPage}

                    \item \textbf{Coursework:} Mathematics, Robotics, Automatic Control Fundamentals, Computer-Based Control, Programming, Computer Vision, Electronics, and Mechanics.
            \end{highlights}
        \end{onecolentry}

    \section{Publications}

        \begin{samepage}
            \begin{twocolentry}{
                April 2024
            }
                \textbf{Exploring the Synergies of Modularization, Interface Standardization, and Service-Orientation in Production System Simulation}

                \vspace{0.10 cm}

                \mbox{Martin Birtic}, \mbox{\textbf{\textit{Pablo Morilla Cabello}}}, \mbox{Ángel Muñoz}, \mbox{Anna Syberfeldt}
            \end{twocolentry}


            \vspace{0.10 cm}

            \begin{onecolentry}
            \textit{DOI: \href{https://doi.org/10.3233/ATDE240164}{10.3233/ATDE240164}}
            \end{onecolentry}
        \end{samepage}

\newpage

    \section{Projects}

        \begin{twocolentry}{


        \textit{\href{https://github.com/PabloMorillaCabello/AGV-Digital-Twin}{github.com/repo}}}
            \textbf{AGV Digital Twin}
            \textit{\href{https://www.linkedin.com/posts/pablomorillacabello_software-opcua-server-activity-7181380999304339457-X081?utm_source=share&utm_medium=member_desktop}{(Project Video)} }
        \end{twocolentry}

        \vspace{0.10 cm}
        \begin{onecolentry}
            \begin{highlights}
                \item This project aims to develop a solution for integrating various software systems using OPC UA standards. It focuses on simulating the functionality of an AGV in a virtual environment, incorporating control capabilities for path planning and collision avoidance while providing a user-friendly remote control interface.
                \item Tools Used: C, JavaScript, MongoDB, Node-RED, RobotStudio, CODESYS, OPC UA.
            \end{highlights}
        \end{onecolentry}

        \vspace{0.2 cm}

        \begin{twocolentry}{
        \textit{\href{https://github.com/PabloMorillaCabello/Mobile-Robotics-Lab}{github.com/repo}}}
            \textbf{Mobile Robotic Lab}
            \textit{\href{https://www.linkedin.com/posts/pablomorillacabello_maerlagauniversity-robot-piero-activity-7160223406435491840-jWpq?utm_source=share&utm_medium=member_desktop}{(Project Video)} }
        \end{twocolentry}

        \vspace{0.10 cm}
        \begin{onecolentry}
            \begin{highlights}
                \item Robot project utilizing a PIERO mobile robot, equipped with sensors and a drive system. Programmed using an Arduino Mega in MATLAB, enabling the implementation of advanced control strategies, including path planning and collision avoidance.
                \item Tools Used: MATLAB, Simulink, Arduino.
            \end{highlights}
        \end{onecolentry}

        \vspace{0.2 cm}

        \begin{twocolentry}{

        \textit{\href{https://github.com/PabloMorillaCabello/Ultra-Adding}{github.com/repo}}}
            \textbf{ULTRA-ADDIN}
            \textit{\href{https://www.linkedin.com/posts/pablomorillacabello_addin-post-git-activity-7135249770603761664-ZB-b?utm_source=share&utm_medium=member_desktop}{(Project Video)} }
        \end{twocolentry}

        \vspace{0.10 cm}
        \begin{onecolentry}
            \begin{highlights}
                \item Development of an add-in for RobotStudio that provides an interface ensuring modularization and interface standardization while keeping the system user-friendly. It supports the research on the synergies between modularization, interface standardization, and service-orientation in production system simulation.
                \item Tools Used: C\#, RobotStudio, Python.
            \end{highlights}
        \end{onecolentry}


        \section{Technologies}

        \begin{onecolentry}
            \textbf{Coding Languages:} C++, C, C\#, Python, Java, JavaScript, RAPID, MATLAB, LaTeX.
        \end{onecolentry}

        \vspace{0.2 cm}

        \begin{onecolentry}
            \textbf{Automation \& Robotics:} RobotStudio, TIA Portal, TwinCAT, CODESYS, Plant Simulation, ROS2, OPC UA, MQTT/HTTP, Node-RED, Simulink, Arduino, Xilinx FPGA.
        \end{onecolentry}

        \vspace{0.2 cm}

        \begin{onecolentry}
            \textbf{AI \& Software:} TensorFlow, PyTorch, OpenCV, Docker, Git/GitHub, MongoDB, Statgraphics, AutoCAD.
        \end{onecolentry}

        \section{Languages}

        \begin{onecolentry}
         Spanish (Native) / English (Fluent) / Swedish (Basic).
        \end{onecolentry}

\end{document}
